\chapter{The near-$126$ reduction and permanent cycle}
\label{chap:near126}

This chapter formalizes Section 7.  Computed Ext groups, products,
differentials, exclusions, and $tmf$ images are explicit evidence inputs.
Logical reductions, representative changes, the Toda argument, the Mahowald
application, and the final contradiction are project proofs.

Every declaration below whose direct dependencies include
Definition~\ref{def:near126-computation-input} is universally quantified over
an explicit parameter $D$ of that type.  Repeated dependencies in one proof
are instantiated with that same $D$; they are not resolved by typeclass search
or a global dataset.  Likewise, every literature or evidence dependency is an
explicit theorem argument.  The final theorem supplies these arguments from
the corresponding projections of its main input.

\section{The inductive criterion and three conditions}

% SOURCE: aimpaper/main.tex:2112-2114, Notation 7.2.
\begin{definition}[Synthetic $\theta_5$ and $\eta$ choices]
  \label{def:synthetic-theta5-choice}
  \uses{def:adams-hj,def:synthetic-adams-ss,def:adams-detection}
  \notready
  A synthetic $\theta_5$ is a class in
  $\pi_{62,64}S^{0,0}$ detected by $h_5^2$; the same symbol denotes each of its
  quotient images.  Let $\eta=[h_1]\in\pi_{1,2}S^{0,0}$.  Every use records
  the chosen representative unless choice independence has already been
  proved.
\end{definition}

% SOURCE: Barratt--Jones--Mahowald inductive method and Burklund--Xu Proposition 7.19; aimpaper/main.tex:2117-2125, whose quantifier is ``for some theta_5''.
\begin{theorem}[BJM/BX criterion for the source choice]
  \label{thm:external-bjm-bx-criterion}
  \uses{def:synthetic-theta5-choice,def:external-result}
  \notready
  An explicit external-result record supplies a distinguished order-two class
  $\theta_5^{\mathrm{src}}$ detected by $h_5^2$ for which:
  \begin{enumerate}
    \item for every integer $r\ge1$, $h_6^2$ survives to $E_{r+3}$ if and
      only if $\lambda\eta(\theta_5^{\mathrm{src}})^2=0$ in
      $\pi_{125,129}(S^{0,0}/\lambda^{r+1})$;
    \item $h_6^2$ is permanent if and only if
      $\lambda\eta(\theta_5^{\mathrm{src}})^2=0$ in
      $\pi_{125,129}S^{0,0}$.
  \end{enumerate}
  The witness and the two equivalences belong to the same source-bearing
  record, which also includes the grading translation from Burklund--Xu to AIM.
\end{theorem}

\begin{proof}
  This node is the located external input, not a project proof.  Retain the
  distinguished class used by the BJM/BX construction, and apply the explicit
  suspension and grading adapter to its finite and infinite statements.  No
  assertion for an arbitrary choice of $\theta_5$ is included here.
\end{proof}

% SOURCE: aimpaper/main.tex:2127-2139, Remarks 7.4-7.5; Xu and IWX inputs.
\begin{theorem}[External $\theta_5$ order and choice comparison]
  \label{thm:external-theta5-order-data}
  \uses{def:synthetic-theta5-choice,def:external-result}
  \notready
  A named Xu/IWX external-result value supplies that every classical
  $\theta_5$ detected by $h_5^2$ has order two and that the difference of any
  two choices has Adams filtration at least six.  The value carries exact
  theorem locators and the comparison from the source's classical convention
  to Definition~\ref{def:synthetic-theta5-choice}.
\end{theorem}

\begin{proof}
  This is a pure literature input.  Store the two quantified conclusions and
  their grading adapter in one named field; do not derive either conclusion
  from the near-$126$ tables.
\end{proof}

% SOURCE: aimpaper/main.tex:2127-2139, Remarks 7.4-7.5; Appendix Tables 3 and 5.
\begin{proposition}[$\theta_5$ order and torsion evidence]
  \label{evidence:theta5-order-torsion}
  \uses{def:near126-computation-input,def:synthetic-theta5-choice,thm:external-theta5-order-data,thm:external-synthetic-einfty-nu,def:external-evidence,data:table-s123,data:table-s124-low}
  \notready
  For a near-$126$ computation input $D$, the named Xu/IWX result and $D$'s
  finite Ext evidence supply: every classical
  $\theta_5$ has order two; the difference of two choices has filtration at
  least six; $\pi_{62,64}S^{0,0}$ has no $\lambda$-torsion; and
  $\pi_{125,130}S^{0,0}$ has no $\lambda$-torsion in the comparison range.
  These facts validate both the order-two hypothesis and the paper's
  one-$\lambda$ shift of the finite BJM/BX statement.
\end{proposition}

\begin{proof}
  Load the located homotopy-group structure theorem and order computation, then
  evaluate the synthetic $E_\infty$ formula on the complete table data in the
  two stated bidegrees.  Exhaustively exclude a class with the required
  torsion length and record the search range in the evidence certificate.
\end{proof}

% SOURCE: aimpaper/main.tex:2133-2139, Remark rem:theta5choice, applying the construction in the proof of Burklund--Xu Proposition 7.19 to an order-two bottom-cell choice.
\begin{theorem}[Total differential identity for an order-two choice]
  \label{thm:external-total-differential-identity}
  \uses{def:synthetic-theta5-choice,def:lambda-rho-delta-triangle,def:external-result}
  \notready
  For every order-two synthetic choice $\theta_5$, the located quadratic-
  construction input supplies the identity
  \[
    \delta_1(h_6^2)=\lambda\eta\theta_5^2
  \]
  in the $\delta_1:S^{0,0}/\lambda\to S^{1,-1}$ total-extension problem, with
  its representative and grading conventions.
\end{theorem}

\begin{proof}
  This is the external total-differential comparison asserted in the cited AIM
  remark: use the chosen order-two class as the bottom-cell restriction in the
  quadratic construction, then translate the target suspension to $S^{1,-1}$.
  The project does not reconstruct that construction internally.
\end{proof}

% SOURCE: aimpaper/main.tex:2133-2139, Remark rem:theta5choice.
\begin{proposition}[Choice independence of the inductive expression]
  \label{prop:theta5-choice-independence}
  \uses{def:near126-computation-input,def:synthetic-theta5-choice,evidence:theta5-order-torsion,thm:external-total-differential-identity}
  \notready
  For a near-$126$ computation input $D$, every synthetic $\theta_5$ has order
  two, and the element
  $\lambda\eta\theta_5^2$ is independent of the choice of $\theta_5$.
  More precisely, any two choices give equal elements in the untruncated
  synthetic sphere and hence equal images in every finite $\lambda$-quotient.
\end{proposition}

\begin{proof}
  Lift the classical order statement through the torsion-free synthetic group.
  The total differential is defined solely by the fixed source class $h_6^2$
  and the connecting map.  Apply the order-two total-differential identity to
  each of two choices; both products equal that same differential.  Naturality
  of the quotient maps gives equality in every finite quotient.
\end{proof}

% SOURCE: aimpaper/main.tex:2117-2139, Theorem thm:bjmbx and Remark rem:theta5choice; project transport from the distinguished source choice.
\begin{theorem}[BJM/BX criterion for any order-two choice]
  \label{thm:bjm-bx-criterion-any-choice}
  \uses{def:near126-computation-input,def:synthetic-theta5-choice,thm:external-bjm-bx-criterion,prop:theta5-choice-independence}
  \notready
  For a near-$126$ computation input $D$ and every order-two synthetic choice
  $\theta_5$:
  \begin{enumerate}
    \item for every integer $r\ge1$, $h_6^2$ survives to $E_{r+3}$ if and
      only if $\lambda\eta\theta_5^2=0$ in
      $\pi_{125,129}(S^{0,0}/\lambda^{r+1})$;
    \item $h_6^2$ is permanent if and only if
      $\lambda\eta\theta_5^2=0$ in $\pi_{125,129}S^{0,0}$.
  \end{enumerate}
\end{theorem}

\begin{proof}
  Choose the distinguished $\theta_5^{\mathrm{src}}$ carried by the external
  criterion.  Fix an integer $r\ge1$.  Choice independence identifies its
  product with the product for the given $\theta_5$, both before and after reduction modulo
  $\lambda^{r+1}$.  Substitute these equalities into the two source
  equivalences.  This transport is a project proof and does not strengthen the
  external source record.
\end{proof}

% SOURCE: aimpaper/main.tex:2181-2204, Proposition prop:possibleh62, statements (3)-(5).
\begin{definition}[The three near-$126$ conditions]
  \label{def:near126-conditions}
  \uses{def:synthetic-theta5-choice,def:adams-detection}
  \notready
  Define:
  \begin{enumerate}
    \item[$C_3$:] $d_6(x_{126,8,4}+x_{126,8})=0$;
    \item[$C_4$:] for some $\theta_5$,
      $\theta_5^2=\lambda^6[h_0^2x_{124,8}]\ne0$ in
      $\pi_{124,128}S^{0,0}$;
    \item[$C_5$:] for some detected $[h_0^2x_{124,8}]$,
      $\lambda^3\eta[h_0^2x_{124,8}]
       =\lambda^6[h_1h_4x_{109,12}]$ in
      $\pi_{125,133}S^{0,0}$.
  \end{enumerate}
  Quantified detected representatives and nonzero assertions are fields of the
  propositions, not informal annotations.
\end{definition}

\section{Load-bearing near-$126$ evidence}

% SOURCE: aimpaper/main.tex:2156-2177, Fact fact:theta5sqAF and following remark; Appendix Tables 5,6,7,9.
\begin{proposition}[Core near-$126$ evidence]
  \label{evidence:near126-core}
  \uses{def:near126-computation-input,def:near126-conditions,data:table-s124-low,data:table-s125-high,data:table-s125-low,data:table-s126-low,thm:appendix-all-rows-encoded,prop:lin-computation-provenance}
  \notready
  For a near-$126$ computation input $D$, its explicit computation evidence
  supplies:
  \begin{enumerate}
    \item $x_{126,8,4}+x_{126,8}$ survives through $E_6$;
    \item $h_1h_4x_{109,12}$ supports no outgoing differential, and the
      exhaustive search leaves exactly two possible incoming killers:
      $d_6(x_{126,8,4}+x_{126,8})$ and $d_{12}(h_6^2)$;
    \item $h_0^2x_{124,8}$ is permanent;
    \item $g^4\Delta h_1g$ is the unique class in its bidegree surviving to
      $E_5$;
    \item $d_3(x_{126,6})$ is one of the two printed nonzero target sums and
      therefore cannot kill $h_1h_4x_{109,12}$.
  \end{enumerate}
  Item (2) deliberately does not call the class permanent before both possible
  incoming differentials are excluded.
\end{proposition}

\begin{proof}
  Run typed queries over the complete four table ranges and their differential
  relation graph.  Check survival, outgoing and incoming ranges separately,
  preserve the finite ambiguity of $d_3(x_{126,6})$, and attach the Zenodo
  disproof identifiers for every excluded potential differential.
\end{proof}

% SOURCE: aimpaper/main.tex:2221-2258, proofs of Lemmas lem:equistate4 and lem:equistate5.
\begin{proposition}[Representative-indeterminacy evidence]
  \label{evidence:near126-indeterminacy}
  \uses{def:near126-computation-input,evidence:theta5-order-torsion,data:table-s124-high,data:table-s124-low,thm:appendix-all-rows-encoded}
  \notready
  For a near-$126$ computation input $D$, its complete stem-$124$ data supply
  that filtration-11 and -12 cycles relevant
  to $C_5$ are killed by $d_2$ or $d_3$, filtration-13 correction terms are
  annihilated by $h_1$, and all remaining corrections begin in filtration at
  least 15 after multiplication by $\lambda^3\eta$.
\end{proposition}

\begin{proof}
  Enumerate every cycle in stem 124 and filtrations 11--13 from both table
  bands, join its outgoing relation and product-by-$h_1$ record, and verify the
  stated alternative for each.  Completeness excludes an unlisted correction.
\end{proof}

% SOURCE: aimpaper/main.tex:2221-2241, Lemma lem:equistate4.
\begin{lemma}[Existential and universal forms of $C_4$]
  \label{lem:near126-equiv-c4}
  \uses{def:near126-computation-input,def:near126-conditions,evidence:theta5-order-torsion}
  \notready
  For a near-$126$ computation input $D$, condition $C_4$ is equivalent to the
  same statement for every choice of $\theta_5$.
\end{lemma}

\begin{proof}
  The difference of two $\theta_5$ choices has filtration at least six, and
  every choice has order two.  Hence their squares differ in filtration at
  least 12.  The displayed detector in $C_4$ has filtration ten and is nonzero,
  so changing $\theta_5$ cannot alter its leading term.
\end{proof}

% SOURCE: aimpaper/main.tex:2244-2258, Lemma lem:equistate5.
\begin{lemma}[Existential and universal forms of $C_5$]
  \label{lem:near126-equiv-c5}
  \uses{def:near126-computation-input,def:near126-conditions,evidence:near126-indeterminacy}
  \notready
  For a near-$126$ computation input $D$, condition $C_5$ is equivalent to the
  statement that its displayed relation holds for every homotopy class
  detected by $h_0^2x_{124,8}$.
\end{lemma}

\begin{proof}
  Two detected representatives differ by a class of filtration at least 11.
  The exhaustive evidence shows that multiplication by $\lambda^3\eta$ sends
  this difference to filtration at least 15.  Since the target leading class
  has filtration 14, the relation is independent of the representative.
\end{proof}

% SOURCE: aimpaper/main.tex:2323-2354, exhaustive theta-square and eta-square analysis; tmf citation.
\begin{proposition}[$\theta_5^2$ and $tmf$ exhaustive evidence]
  \label{evidence:theta5-square-tmf}
  \uses{def:near126-computation-input,evidence:near126-core,evidence:near126-indeterminacy,data:table-s125-low,data:table-s125-high,def:external-result,def:external-evidence}
  \notready
  For a near-$126$ computation input $D$, its explicit table and $tmf$ inputs
  supply:
  \begin{enumerate}
    \item the only possibilities for $\theta_5^2$ are
      $\lambda^6[h_0^2x_{124,8}]$,
      $\lambda^9[e_0\Delta h_6g]$, or a $\lambda^{10}$-multiple;
    \item $h_1e_0\Delta h_6g=0$;
    \item the only relevant high-filtration $\lambda^{10}$-multiple is the
      class represented by $\lambda^{20}g^4\Delta h_1g$;
    \item its $tmf$ Hurewicz image is nonzero, while every $\theta_5$ maps to
      zero in $tmf$;
    \item if $C_3$ fails, the unique alternative incoming differential is
      $d_6(x_{126,8,4}+x_{126,8})=h_1h_4x_{109,12}$.
  \end{enumerate}
\end{proposition}

\begin{proof}
  Query the complete filtration ranges and product table to produce the three
  alternatives and unique surviving high class.  Attach independent, located
  $tmf$ Hurewicz evidence for its image and the image of $\theta_5$.  Finally
  use the incoming-differential completeness certificate for the stated
  alternative.  Each negative conclusion retains its finite search bounds.
\end{proof}

% SOURCE: aimpaper/main.tex:2181-2204 and 2300-2354, Proposition prop:possibleh62.
\begin{proposition}[Only possible differential on $h_6^2$]
  \label{prop:possible-h62}
  \uses{def:near126-computation-input,def:near126-conditions,thm:bjm-bx-criterion-any-choice,thm:external-total-differential-identity,thm:delta-classical-equivalence,evidence:near126-core,evidence:theta5-square-tmf,lem:near126-equiv-c4,lem:near126-equiv-c5}
  \notready
  For a near-$126$ computation input $D$, exactly one of the following holds:
  \begin{enumerate}
    \item $h_6^2$ is a permanent cycle;
    \item $d_{12}(h_6^2)=h_1h_4x_{109,12}$ is nonzero.
  \end{enumerate}
  Moreover, the second statement holds if and only if
  $C_3$, $C_4$, and $C_5$ all hold.
\end{proposition}

\begin{proof}
  Identify the total $\delta_1$ differential with
  $\lambda\eta\theta_5^2$ and then with the corresponding classical
  differential coset.  Under $C_4$ and $C_5$, its target is the stated
  filtration-14 class; $C_3$ and the two-killer certificate make its
  $\lambda^{10}$ multiple nonzero, giving $d_{12}$.  Conversely, use the
  exhaustive three-way classification of $\theta_5^2$.  Failure of $C_4$ or
  $C_5$ forces the total expression into the unique $tmf$-detected high class,
  where its zero $tmf$ image makes it zero.  Failure of $C_3$ activates the
  unique $d_6$ alternative and again kills the relevant synthetic multiple.
  The BJM/BX criterion then gives permanence.  The two alternatives are
  exclusive because a source supporting a nonzero $d_{12}$ is not permanent.
\end{proof}

\section{Reduction of $C_5$ to a two-extension}

% SOURCE: aimpaper/main.tex:2377-2385, Fact fact:x1239; Tables S123 and S125.19.
\begin{proposition}[$x_{123,9}$ computation evidence]
  \label{evidence:x1239}
  \uses{def:near126-computation-input,data:table-s123,data:table-s125-low,thm:appendix-all-rows-encoded}
  \notready
  For a near-$126$ computation input $D$, the class
  $x_{123,9}+h_0x_{123,8}$ survives through $E_{12}$ and has no incoming
  differential in $D$, and
  \[
    d_2(x_{125,8})=
      h_1(x_{123,9}+h_0x_{123,8})+h_0^2x_{124,8}.
  \]
\end{proposition}

\begin{proof}
  Retrieve the named rows and shared differential relation, then use complete
  source ranges to certify the negative incoming claim through the relevant
  maximum length.
\end{proof}

% SOURCE: aimpaper/main.tex:2387-2446, Lemma lem:x1239.
\begin{lemma}[The classes $\alpha_1,\alpha_2,\alpha_3$]
  \label{lem:near126-alpha-relations}
  \uses{def:near126-computation-input,evidence:x1239,thm:external-synthetic-einfty-quotient,evidence:near126-indeterminacy,data:table-s123,data:table-s124-high}
  \notready
  For a near-$126$ computation input $D$, there exists
  $\alpha_1=[x_{123,9}+h_0x_{123,8}]$ in
  $\pi_{123,132}(S^{0,0}/\lambda^9)$ such that for every
  $[h_0^2x_{124,8}]$ there are
  $\alpha_2\in\pi_{124,137}(S^{0,0}/\lambda^9)$ and
  $\alpha_3\in\pi_{125,140}(S^{0,0}/\lambda^9)$ satisfying
  \[
  \begin{aligned}
    \lambda^3\eta\alpha_1
      &=\lambda^3[h_0^2x_{124,8}]+\lambda^6\alpha_2,\\
    \eta\alpha_2&=\lambda\alpha_3,
  \end{aligned}
  \]
  and $\lambda^3\alpha_1[h_0]=0$.  Moreover
  $\operatorname{AF}(\alpha_3)>14$, so
  $\lambda^7\alpha_3$ lies strictly above filtration $14$ in the target
  group.
\end{lemma}

\begin{proof}
  Lift $\alpha_1$ first modulo $\lambda^{11}$ using its $E_{12}$ survival.
  The computed $d_2$ identifies the leading $h_1$ product; enumerate the
  remaining filtration-13 correction and choose $\alpha_2$, whose $h_1$
  product is zero in associated-graded filtration $14$ and is
  $\lambda$-divisible, to obtain $\alpha_3$ in filtration at least $15$.
  For the final zero,
  evaluate the finite-quotient $E_\infty$ formula and exclude the two lower
  candidates using the recorded $d_7$ and $d_3$ relations.  What remains is
  $\lambda^{10}$-divisible modulo $\lambda^{11}$ and hence maps to zero modulo
  $\lambda^9$.
\end{proof}

% SOURCE: aimpaper/main.tex:2466-2468, Fact fact:h02x1259; Table S125.19.
\begin{proposition}[$h_0^2x_{125,9,2}$ evidence]
  \label{evidence:h02x1259}
  \uses{def:near126-computation-input,data:table-s125-low,thm:appendix-all-rows-encoded}
  \notready
  For a near-$126$ computation input $D$, the class $h_0^2x_{125,9,2}$
  survives through $E_5$ and no incoming
  differential hits it.  Its printed $d_5?$ status remains unresolved and is
  not asserted nonzero.
\end{proposition}

\begin{proof}
  Read the typed row and query every possible incoming source in the complete
  adjacent-stem ranges.  Preserve the unresolved-outgoing constructor for the
  $d_5$ column.
\end{proof}

% SOURCE: noncircular repair of aimpaper/main.tex:2493-2496 using lines 2427-2429 and 2579.
\begin{lemma}[Leading-term nonvanishing for the Toda product]
  \label{lem:toda-h0-leading-nonvanishing}
  \uses{def:near126-computation-input,def:near126-conditions,lem:near126-equiv-c5,lem:near126-alpha-relations,evidence:near126-core,thm:external-synthetic-einfty-quotient,thm:toda-product-identities}
  \notready
  For a near-$126$ computation input $D$, assume $C_3$ and the universal form
  of $C_5$.  For the classes from
  Lemma~\ref{lem:near126-alpha-relations}, the element
  \[
    \beta=\lambda^6[h_1h_4x_{109,12}]+\lambda^7\alpha_3
      \in\pi_{125,133}(S^{0,0}/\lambda^9)
  \]
  is nonzero and has filtration-$14$ leading term
  $\lambda^6h_1h_4x_{109,12}$.  Every element of the defined bracket
  $\langle\lambda^3\alpha_1,[h_0],\eta\rangle$ has $[h_0]$-product
  $\beta$; consequently the bracket does not contain zero.
\end{lemma}

\begin{proof}
  Multiply the first $\alpha$ relation by $\eta$, replace
  $\eta\alpha_2$ by $\lambda\alpha_3$, and apply the universal $C_5$
  relation.  The computation evidence leaves only a $d_6$ or $d_{12}$ as an
  incoming killer of $h_1h_4x_{109,12}$.  Condition $C_3$ removes the former,
  while the finite $\lambda^9$ quotient does not see the latter; hence the
  first summand is nonzero in filtration $14$.  The $\alpha$ lemma places the
  second summand in filtration at least $15$, so separatedness prevents
  cancellation.  Finally apply the stated Toda shuffle and
  $\eta^2$ identities.  This argument uses neither the unresolved $d_5$ nor
  the Toda-detection conclusion below.
\end{proof}

% SOURCE: aimpaper/main.tex:2470-2567, Lemma lem:toda2ext; Moss Theorem 1.2 and listed Ext products.
\begin{lemma}[Toda bracket detects the two-extension source]
  \label{lem:toda-two-extension}
  \uses{def:near126-computation-input,def:near126-conditions,lem:near126-equiv-c5,lem:near126-alpha-relations,lem:toda-h0-leading-nonvanishing,thm:moss-convergence-adapter,thm:toda-product-identities,evidence:near126-core,evidence:h02x1259,def:external-evidence}
  \notready
  For a near-$126$ computation input $D$, assuming $C_3$ and the universal
  form of $C_5$, the synthetic Toda bracket
  \[
    \langle\lambda^3\alpha_1,[h_0],\eta\rangle
       \subseteq\pi_{125,132}(S^{0,0}/\lambda^9)
  \]
  is defined, does not contain zero, and is detected by
  $\lambda^4h_0^2x_{125,9,2}$.
\end{lemma}

\begin{proof}
  The zero product from the $\alpha$ lemma defines the bracket, and the
  leading-term lemma proves its $[h_0]$ product is nonzero without assuming the
  unresolved $d_5$.  Exhaustively enumerate the filtration-at-most-12
  candidates.  Exclude the filtration-7 class by its $d_3$ torsion length.
  For the filtration-9 class, apply the located Moss theorem to the displayed
  Massey product, control Toda indeterminacy, and use the explicit Ext product
  inequality to contradict $C_5$.  The sole remaining candidate is the stated
  filtration-11 class.
\end{proof}

% SOURCE: aimpaper/main.tex:2571-2580, proof of Corollary cor:2ext125, especially the comparison of all representatives detected by lambda^4 h_0^2 x_{125,9,2}; Appendix Table S125.19.
\begin{lemma}[Indeterminacy of the two-extension source]
  \label{lem:near126-two-extension-indeterminacy}
  \uses{def:near126-computation-input,def:near126-conditions,lem:near126-equiv-c5,lem:toda-h0-leading-nonvanishing,lem:toda-two-extension,evidence:h02x1259,evidence:near126-core,data:table-s125-low,thm:appendix-all-rows-encoded,thm:external-synthetic-einfty-quotient,def:external-evidence}
  \notready
  For a near-$126$ computation input $D$, assume $C_3$ and the universal form
  of $C_5$.  Let $u_0$ be a bracket
  representative supplied by Lemma~\ref{lem:toda-two-extension}, and let $u$
  be any class in $\pi_{125,132}(S^{0,0}/\lambda^9)$ detected by
  $\lambda^4h_0^2x_{125,9,2}$.  If

  \[
    u_0[h_0]=\lambda^6[v_0]
  \]
  for a representative $[v_0]$ detected by $h_1h_4x_{109,12}$, then there is
  a representative $[v_u]$ detected by the same Ext class such that

  \[
    u[h_0]=\lambda^6[v_u].
  \]
  The filtration-$14$ leading term is nonzero and is unchanged by replacing
  $u_0$ with $u$.
\end{lemma}

\begin{proof}
  The difference $u-u_0$ is detected in strictly higher Adams filtration than
  $\lambda^4h_0^2x_{125,9,2}$.  Enumerate, from the complete stem-$125$ data,
  every class that can occur in this detection indeterminacy and evaluate its
  $[h_0]$ product using $D.\mathrm{twoExtensionIndeterminacyProducts}$, the
  named finite Ext-product record.  None has a
  filtration-$14$ component detected by
  $\lambda^6h_1h_4x_{109,12}$; every remaining product is precisely a change
  of representative for the target detector.  Absorb their sum into $[v_0]$
  to obtain $[v_u]$.  The leading-term lemma proves that the common
  filtration-$14$ component is nonzero, so it cannot be cancelled by this
  change.  The finite enumeration and product records are external evidence;
  the indeterminacy deduction is the project proof recorded here.
\end{proof}

% SOURCE: aimpaper/main.tex:2478-2480, Remark rem:h02x1259, now deduced after Lemma lem:toda2ext.
\begin{corollary}[Conditional vanishing of the unresolved $d_5$]
  \label{cor:h02x1259-d5-zero}
  \uses{def:near126-computation-input,lem:near126-equiv-c5,lem:toda-two-extension,evidence:h02x1259,thm:external-synthetic-einfty-quotient}
  \notready
  For a near-$126$ computation input $D$, under $C_3$ and $C_5$, the class
  $\lambda^4h_0^2x_{125,9,2}$ detects a nonzero class modulo $\lambda^9$.
  Equivalently, the unresolved classical differential satisfies
  $d_5(h_0^2x_{125,9,2})=0$ under these hypotheses.
\end{corollary}

\begin{proof}
  First use Lemma~\ref{lem:near126-equiv-c5} to transport the assumed
  existential $C_5$ to the universal form required by the Toda lemma.  That
  lemma supplies a nonzero bracket element detected by the displayed
  class.  Apply the finite-quotient $E_\infty$ formula: because the underlying
  row survives through $E_5$ and has no incoming differential, this
  $\lambda^4$ multiple exists modulo $\lambda^9$ exactly when its possible
  outgoing $d_5$ vanishes.
\end{proof}

% SOURCE: aimpaper/main.tex:2571-2581, Corollary cor:2ext125.
\begin{corollary}[The required two-extension]
  \label{cor:near126-two-extension}
  \uses{def:near126-computation-input,lem:near126-equiv-c5,lem:toda-two-extension,lem:near126-two-extension-indeterminacy,cor:h02x1259-d5-zero}
  \notready
  For a near-$126$ computation input $D$, under $C_3$ and $C_5$, every class
  detected by
  $\lambda^4h_0^2x_{125,9,2}$ satisfies
  \[
    [\lambda^4h_0^2x_{125,9,2}][h_0]
      =\lambda^6[h_1h_4x_{109,12}]\ne0
  \]
  in $\pi_{125,133}(S^{0,0}/\lambda^9)$, for a suitable detected target
  representative.
\end{corollary}

\begin{proof}
  First transport $C_5$ to its universal form with
  Lemma~\ref{lem:near126-equiv-c5}.  Choose the bracket representative supplied
  by the Toda lemma and compute its
  $[h_0]$ product.  Absorb the higher-filtration $\lambda^7\alpha_3$ term by
  changing the detected target representative.  Apply the source-indeterminacy
  lemma to transport this relation from that bracket representative to every
  class detected by $\lambda^4h_0^2x_{125,9,2}$.  Its preserved
  filtration-$14$ leading term gives the asserted nonvanishing.
\end{proof}

\section{The $\nu_{\mathrm{Hopf}}$ extension and contradiction}

% SOURCE: aimpaper/main.tex:2613-2615, Fact fact:h1x1217; Table S122.
\begin{proposition}[$h_1x_{121,7}$ evidence]
  \label{evidence:h1x1217}
  \uses{def:near126-computation-input,data:table-s122,thm:appendix-all-rows-encoded}
  \notready
  For a near-$126$ computation input $D$, the class $h_1x_{121,7}$ survives
  through $E_6$ and has no incoming
  differential in the complete range.
\end{proposition}

\begin{proof}
  Query its table row and all degree-compatible incoming relation identifiers;
  use table completeness to certify the negative result.
\end{proof}

% SOURCE: aimpaper/main.tex:2617-2673, missing obstruction check in Lemma lem:nuext125.
\begin{proposition}[Obstruction vanishing for the mod-$\lambda^5$ lift]
  \label{evidence:nu-extension-obstruction-vanishing}
  \uses{def:near126-computation-input,data:table-s125-low,data:table-s126-low,data:table-cnu126,thm:appendix-all-rows-encoded}
  \notready
  For a near-$126$ computation input $D$ and the relation produced modulo
  $\lambda^3$ by the $S^0/\nu_{\mathrm{Hopf}}$
  $d_3$, every filtration-10 correction that could obstruct a lift modulo
  $\lambda^5$ either has no detected homotopy class or can be killed by changing
  the lift.  Therefore a lift exists with no additional $[\lambda x]$ term.
\end{proposition}

\begin{proof}
  Enumerate the complete filtration-10 source group, compute the relevant
  $\rho$ obstruction and each candidate's synthetic survival from the table
  relations, and give an explicit correction or vanishing certificate for
  every element.  This evidence is independent of the formal properties of
  $\rho$; it fills the unsupported ``upon inspection'' step at
  \texttt{aimpaper/main.tex:2661--2665}.
\end{proof}

% SOURCE: aimpaper/main.tex:2617-2673, Lemma lem:nuext125; computed cofiber d_3 and no-crossing evidence.
\begin{lemma}[The required Hopf extension]
  \label{lem:near126-nu-extension}
  \uses{def:near126-computation-input,evidence:h1x1217,thm:generalized-mahowald,lem:normalized-hopf-exactness-case,def:hopf-cofiber-cell-notation,data:table-cnu126,evidence:nu-extension-obstruction-vanishing,prop:synthetic-einfty-map-compatibility}
  \notready
  For a near-$126$ computation input $D$, there is a class
  $[\lambda^4h_1x_{121,7}]\in
  \pi_{122,126}(S^{0,0}/\lambda^9)$ such that
  \[
    [\lambda^4h_1x_{121,7}][h_2]
       =\lambda[\lambda^5h_0^2x_{125,9,2}]
  \]
  in $\pi_{125,130}(S^{0,0}/\lambda^9)$.
\end{lemma}

\begin{proof}
  Instantiate the Hopf cofiber triangle with the zero-tagged exactness datum
  of Lemma~\ref{lem:normalized-hopf-exactness-case}, so the Mahowald package
  has $(e(f),e(g),e(h))=(1,0,0)$ and uses that one lifted triangle.  Apply the
  generalized Mahowald trick to this triangle and the
  typed $S^0/\nu_{\mathrm{Hopf}}$ differential
  displayed in \texttt{aimpaper/main.tex:2638--2648}.  Its edge extensions and
  no-crossing certificate give an $([h_2],E_4)$ extension, hence a relation
  modulo $\lambda^3$.  Use the independent obstruction-vanishing certificate
  to lift it modulo $\lambda^5$, then multiply by $\lambda^4$ and apply quotient
  compatibility to obtain the displayed relation modulo $\lambda^9$.
\end{proof}

% SOURCE: aimpaper/main.tex:2689-2696, Fact fact:stem122; Table S122.
\begin{proposition}[Stem-$122$ permanent-class evidence]
  \label{evidence:stem122-permanent}
  \uses{def:near126-computation-input,data:table-s122,thm:appendix-all-rows-encoded}
  \notready
  For a near-$126$ computation input $D$, the classes $h_6Md_0$ and
  $h_5x_{91,11}$ are permanent in the classical Adams spectral sequence
  encoded by $D$.
\end{proposition}

\begin{proof}
  Retrieve their rows, certify no outgoing or incoming differential through
  the complete source ranges, and validate the permanent status against the
  program artifact identifiers.
\end{proof}

% SOURCE: aimpaper/main.tex:2707-2728, two ``upon inspection'' steps in the proof of Proposition prop:state5false.
\begin{proposition}[Stem-$122$ candidate and product exhaustion]
  \label{evidence:stem122-product-exhaustion}
  \uses{def:near126-computation-input,data:table-s122,data:table-s125-low,evidence:stem122-permanent,thm:appendix-all-rows-encoded,def:external-evidence}
  \notready
  For a near-$126$ computation input $D$, in the bidegree and
  $\lambda$-quotient occurring in the proof of
  Proposition \texttt{prop:state5false}, every nonzero representative of
  $[\lambda^4h_1x_{121,7}][h_0]$ of Adams filtration at most $12$ is exactly
  one of
  \[
    \lambda^6[h_6Md_0]\quad\text{or}\quad
    \lambda^7[h_5x_{91,11}].
  \]
  The typed product and differential evidence further shows, in either case,
  that $\lambda^4[h_1h_4x_{109,12}]$ is a
  $\lambda[h_2]$-multiple in the untruncated synthetic sphere.  This record
  includes the filtration-$13$ candidate search and the $d_2/d_4$ exclusions
  used in lines 2716--2728.
\end{proposition}

\begin{proof}
  Enumerate the complete stem-$122$ permanent representatives through
  filtration $12$, evaluate the typed $h_0$ and $h_2$ products, and transport
  each known differential through the shared relation identifiers.  Query the
  complete stem-$125$ ranges for filtration-$13$ corrections and prove that
  each is killed by the recorded $d_2$ or $d_4$.  The resulting finite
  certificate has exactly the two source cases and the asserted common
  divisibility conclusion.
\end{proof}

% SOURCE: aimpaper/main.tex:1398-1405, Example exam:extonEn(2), for the
% identity \widehat{\nu_{\mathrm{Hopf}}}=[h_2] and its grading;
% aimpaper/main.tex:929-961, Notation not:fhat, for normalization.
\begin{proposition}[Normalized Hopf-map detection]
  \label{evidence:normalized-hopf-detection}
  \uses{def:synthetic-hat-map,def:adams-detection,thm:external-synthetic-triangle-lift,def:cofiber-triangle,def:external-evidence}
  \notready
  Explicit detection evidence retains one chosen application of the synthetic
  triangle-lift theorem to the appropriate rotation of the classical Hopf
  cofiber triangle.  It identifies that stored triangle's Hopf component with
  \[
    \widehat{\nu_{\mathrm{Hopf}}}=[h_2]:S^{3,4}\longrightarrow S^{0,0},
  \]
  where $[h_2]$ is detected by $h_2$ in the source's stated synthetic
  tridegree.  With this fixed normalization,
  $\nu(\nu_{\mathrm{Hopf}})=\lambda[h_2]$.
\end{proposition}

\begin{proof}
  Use the explicit normalization and tridegree calculation in Example
  \texttt{exam:extonEn(2)} at \texttt{aimpaper/main.tex:1398--1405}.  Store
  the entire chosen lifted-triangle witness, the detected class, the equality
  with its Hopf component, and the factorization by $\lambda$ as one
  source-bearing evidence value.
\end{proof}

% SOURCE: aimpaper/main.tex:1398-1405 and 929-961; project packaging of the
% zero-homology rotation needed by the normalized-cofiber interface.
\begin{lemma}[Normalized Hopf cofiber exactness case]
  \label{lem:normalized-hopf-exactness-case}
  \uses{evidence:normalized-hopf-detection,def:normalized-synthetic-exactness-case,prop:positive-adams-filtration-homology-zero,def:cofiber-triangle,def:hf2-homology,thm:external-synthetic-triangle-lift}
  \notready
  The classical Hopf cofiber triangle
  \[
    S^3\xrightarrow{\nu_{\mathrm{Hopf}}}S^0\xrightarrow g
    S^0/\nu_{\mathrm{Hopf}}\xrightarrow hS^4
  \]
  carries the zero constructor of
  Definition~\ref{def:normalized-synthetic-exactness-case}.  Explicitly,
  \[
    H_*\nu_{\mathrm{Hopf}}=0,
    \qquad
    (e(\nu_{\mathrm{Hopf}}),e(g),e(h))=(1,0,0),
  \]
  and its long exact homology sequence supplies the rotated short exact
  sequence.  The constructor's lifted-triangle field is exactly the witness
  retained by Proposition~\ref{evidence:normalized-hopf-detection}, and its
  normalized Hopf component is therefore equal to $[h_2]$.
\end{lemma}

\begin{proof}
  The filtration-one detection by $h_2$ gives
  $e(\nu_{\mathrm{Hopf}})=1$, while degree considerations make the Hopf map
  zero on $H\mathbb F_2$-homology.  In the cofiber long exact sequence, $g$
  is the nonzero bottom-cell inclusion on homology and $h$ is the nonzero
  top-cell quotient.  The contrapositive of
  Proposition~\ref{prop:positive-adams-filtration-homology-zero} therefore
  gives Adams filtration zero and $e(g)=e(h)=0$.  The remaining
  long-exact-sequence terms split by degree into
  the required short exact sequence.  Package these facts and the displayed
  $e$-pattern around the precise lifted-triangle witness retained by the
  normalized-detection evidence; do not invoke the nonunique lift theorem a
  second time.  The evidence's component equality becomes the zero tag's
  equality $\widehat{\nu_{\mathrm{Hopf}}}=[h_2]$.
\end{proof}

% SOURCE: aimpaper/main.tex:1398-1405 and 929-961; the erroneous cofiber
% conclusion at aimpaper/main.tex:2730-2732 is repaired: the cofiber equivalent
% to nu(S^0/nu) is C([h_2]), not C(lambda[h_2]).
\begin{proposition}[Normalized synthetic Hopf cofiber]
  \label{prop:synthetic-hopf-cofiber-equivalence}
  \uses{evidence:normalized-hopf-detection,lem:normalized-hopf-exactness-case,prop:synthetic-hat-cofiber,def:cofiber-triangle,def:hopf-cofiber-cell-notation}
  \notready
  With the explicit detection identity
  $\widehat{\nu_{\mathrm{Hopf}}}=[h_2]$ and
  $\nu(\nu_{\mathrm{Hopf}})=\lambda[h_2]$, the normalized cofiber satisfies
  \[
    C([h_2])\simeq\nu(S^0/\nu_{\mathrm{Hopf}}).
  \]
  The factorization $\lambda[h_2]=[h_2]\circ\lambda$ induces a comparison map
  $C(\lambda[h_2])\to C([h_2])$; it does not assert that these two cofibers are
  equivalent.  Consequently, an element divisible by $\lambda[h_2]$ maps to
  zero in $C([h_2])$ after absorbing its factor of $\lambda$ into the source of
  $[h_2]$.
\end{proposition}

\begin{proof}
  Apply Proposition~\ref{evidence:normalized-hopf-detection} to identify the
  chosen normalized lift of the filtration-one Hopf map with $[h_2]$.
  Lemma~\ref{lem:normalized-hopf-exactness-case} supplies the zero-tagged
  exactness datum, including $e(g)=0$, the exact lifted triangle retained by
  the detection evidence, and its component equality with $[h_2]$.  Applying that proposition
  to this datum identifies $C([h_2])$, with no weight shift, with the synthetic
  analogue of the classical Hopf cofiber.  The equality
  $\nu(\nu_{\mathrm{Hopf}})=\lambda[h_2]$ gives the displayed factorization.
  Functoriality of cofibers gives the comparison map, while direct substitution
  shows that every $\lambda[h_2]$-multiple is already an $[h_2]$-multiple.
\end{proof}

% SOURCE: aimpaper/main.tex:2698-2778 and Table Table:Cnu126.
\begin{proposition}[$S^0/\nu$ short-incoming exclusion]
  \label{evidence:cnu126-short-incoming-exclusion}
  \uses{def:near126-computation-input,data:table-cnu126,thm:appendix-all-rows-encoded,prop:lin-computation-provenance}
  \notready
  For a near-$126$ computation input $D$, in its classical Adams spectral
  sequence of $S^0/\nu_{\mathrm{Hopf}}$, the class
  $h_1h_4x_{109,12}[0]$ is not hit by any $d_r$ with $r\le5$.  This negative
  statement is a range-completeness query over all 26 table rows, not a claim
  read from a single row.
\end{proposition}

\begin{proof}
  Enumerate every degree-compatible source in filtrations 9--14, follow its
  relation status, and prove none has target the named bottom-cell class with
  length at most five.  Table and artifact completeness make the finite
  negative query exhaustive.
\end{proof}

% SOURCE: aimpaper/main.tex:2208-2210 and 2698-2778, Proposition prop:state5false.
\begin{proposition}[$C_3$ excludes $C_5$]
  \label{prop:near126-c3-not-c5}
  \uses{def:near126-computation-input,def:near126-conditions,lem:near126-equiv-c5,cor:near126-two-extension,lem:near126-nu-extension,evidence:near126-core,evidence:stem122-product-exhaustion,prop:synthetic-hopf-cofiber-equivalence,thm:external-synthetic-rigidity,evidence:cnu126-short-incoming-exclusion,def:external-evidence}
  \notready
  For a near-$126$ computation input $D$, if $C_3$ holds, then $C_5$ is false.
\end{proposition}

\begin{proof}
  Assume $C_3$ and the universal form of $C_5$.  Compose the $[h_2]$ relation
  with the $[h_0]$ extension to obtain the nonzero
  $\lambda^8[h_1h_4x_{109,12}]$.  Exhaustive stem-$122$ evidence says the
  intermediate $[h_0]$ multiple is represented only by
  $\lambda^6[h_6Md_0]$ or $\lambda^7[h_5x_{91,11}]$.  Product and filtration
  evidence in either case makes
  $\lambda^4[h_1h_4x_{109,12}]$ a $\lambda[h_2]$ multiple in the untruncated
  sphere.  It is therefore also an $[h_2]$-multiple and maps to zero in
  $C([h_2])\simeq\nu(S^0/\nu_{\mathrm{Hopf}})$.  Rigidity would then require an
  incoming differential of length at most five on the corresponding
  bottom-cell class in $S^0/\nu_{\mathrm{Hopf}}$.  The complete Table 1 exclusion
  contradicts this.  Hence $C_5$ cannot hold.
\end{proof}

% SOURCE: aimpaper/main.tex:128-130 and 2106-2214, Theorems thm:h62 and thm:126survives.
\begin{theorem}[$h_6^2$ is a permanent cycle]
  \label{thm:h6-square-permanent}
  \uses{def:kervaire-main-input,prop:possible-h62,prop:near126-c3-not-c5}
  \notready
  For every main input $I$, the class $h_6^2$ is a permanent cycle in
  $I.\mathrm{classicalSphere}$, the classical mod--$2$ Adams spectral sequence
  of the sphere.
\end{theorem}

\begin{proof}
  Put $D=I.\mathrm{computation}$.  Instantiate every literature leaf with its
  named projection from $I.\mathrm{literature}$, every table and finite-search
  theorem with this same $D$, and every spectral-sequence comparison through
  $I$'s coherence fields.  If the nonzero
  $d_{12}$ alternative of
  Proposition~\ref{prop:possible-h62} held, then $C_3$, $C_4$, and $C_5$ would
  all hold.  Proposition~\ref{prop:near126-c3-not-c5} contradicts the first and
  third simultaneously.  Therefore the differential alternative is false,
  and the exclusive permanent-cycle alternative holds.
\end{proof}
